\section{Ladders of levels}\label{sec:bounded-interface}

A ladder is a sequence of levels, each adding a distinction that the previous ones lacked.
This section proves the two theorems that bound ladders built from a fixed packaging map or a fixed interface.
Both are elementary; their content lies in what they force an unbounded ladder to change.

\subsection{Idempotent packaging saturates}

\begin{lemma}\label{lem:idempotent-iterates}
If $e\circ e=e$, then $e^{\circ n}=e$ for all $n\ge 1$.
\end{lemma}

\begin{proof}
By induction on $n$; see \Cref{app:packdef}.
\end{proof}

\begin{proof}[Proof of \Cref{thm:NG_LADDER_IDEM}]
This is \Cref{lem:idempotent-iterates}.
In particular, any quantity that depends on a state only through its packaged value $e^{\circ n}(s)$ takes the same value for all $n\ge 1$.
\end{proof}

\subsection{A fixed interface has finitely many definable predicates}

\begin{lemma}\label{lem:definability-bijection}
Let $f:X\to Y$ be any map. The assignment $A\mapsto\1_A\circ f$ is a bijection from the subsets of $\im(f)$ onto $\Def(f)$. It is an order isomorphism: $\1_A\circ f\le\1_B\circ f$ pointwise if and only if $A\subseteq B$.
\end{lemma}

\begin{proof}
See \Cref{app:packdef}. Only the values of $g$ on the image of $f$ affect $g\circ f$, and every point of the image is hit by some $x$.
\end{proof}

\begin{proof}[Proof of \Cref{thm:NG_LADDER_BOUNDED_INTERFACE}]
By \Cref{lem:definability-bijection}, $|\Def(f)|$ is the number of subsets of $\im(f)$, which is $2^{|\im(f)|}$.
An infinite sequence of pairwise distinct elements would make the finite set $\Def(f)$ infinite, which is impossible.
A strictly increasing chain consists of pairwise distinct elements, so none is infinite either.
\end{proof}

\Cref{fig:definability} illustrates the bijection.
The order isomorphism gives a sharper bound on chains.

\begin{figure}[t]
\centering
\begin{tikzpicture}[
  pt/.style={circle,fill=black,inner sep=1.3pt},
  fib/.style={rounded corners=5pt,draw,thick},
  ys/.style={circle,draw,thick,minimum size=7mm,inner sep=0pt,font=\small},
  lens/.style={-{Stealth[length=2mm]},densely dashed,gray!70!black},
  lbl/.style={font=\footnotesize}
]
\begin{scope}
\node[fib,fill=green!12,minimum width=2.2cm,minimum height=1.25cm] (F1) at (0,0) {};
\node[fib,fill=green!12,minimum width=1.4cm,minimum height=1.25cm] (F2) at (2.5,0) {};
\node[fib,fill=white,minimum width=1.9cm,minimum height=1.25cm] (F3) at (4.75,0) {};
\foreach \x/\y in {-0.7/0.2,-0.2/-0.25,0.35/0.25,0.75/-0.2} \node[pt] at (\x,\y) {};
\foreach \x/\y in {2.25/0.2,2.75/-0.2} \node[pt] at (\x,\y) {};
\foreach \x/\y in {4.2/-0.15,4.7/0.25,5.25/-0.2} \node[pt] at (\x,\y) {};
\node[lbl,anchor=east] at (-1.3,0) {$X$};
\node[ys,fill=green!12] (y1) at (0,-2.0) {$y_1$};
\node[ys,fill=green!12] (y2) at (2.5,-2.0) {$y_2$};
\node[ys] (y3) at (4.75,-2.0) {$y_3$};
\node[ys] (y4) at (6.6,-2.0) {$y_4$};
\node[lbl,anchor=east] at (-1.3,-2.0) {$Y$};
\draw[lens] (F1) -- (y1);
\draw[lens] (F2) -- (y2);
\draw[lens] (F3) -- (y3);
\node[lbl,align=center] at (6.6,-1.05) {not in $\im(f)$};
\node[lbl] at (2.4,-2.85) {$A=\{y_1,y_2\}$ \ gives the predicate $\1_A\circ f$ (shaded fibers)};
\end{scope}
\begin{scope}[xshift=10.6cm,yshift=-1.0cm,lbl/.style={font=\scriptsize,fill=white,inner sep=1.5pt}]
\node[lbl] (e) at (0,-1.7) {$\emptyset$};
\node[lbl] (a1) at (-1.45,-0.55) {$\{y_1\}$};
\node[lbl] (a2) at (0,-0.55) {$\{y_2\}$};
\node[lbl] (a3) at (1.45,-0.55) {$\{y_3\}$};
\node[lbl] (b12) at (-1.45,0.6) {$\{y_1,y_2\}$};
\node[lbl] (b13) at (0,0.6) {$\{y_1,y_3\}$};
\node[lbl] (b23) at (1.45,0.6) {$\{y_2,y_3\}$};
\node[lbl] (t) at (0,1.75) {$\im(f)$};
\begin{scope}[on background layer]
\foreach \u/\v in {e/a1,e/a2,e/a3,a1/b12,a1/b13,a2/b12,a2/b23,a3/b13,a3/b23,b12/t,b13/t,b23/t} \draw[gray!60] (\u) -- (\v);
\draw[very thick,green!50!black] (e) -- (a1) -- (b12) -- (t);
\end{scope}
\node[font=\footnotesize,align=center] at (0,-2.35) {$2^3=8$ definable predicates};
\end{scope}
\end{tikzpicture}
\caption{Definable predicates of a lens with $|\im(f)|=3$. Left: a predicate definable through $f$ is constant on fibers, so it is determined by the set $A\subseteq\im(f)$ of images on which it equals one; values of $g$ off the image, such as at $y_4$, are irrelevant. Right: the $8$ definable predicates ordered pointwise form the lattice of subsets of $\im(f)$. A strictly increasing chain, such as the highlighted one, has at most $|\im(f)|+1=4$ elements.}\label{fig:definability}
\end{figure}

\begin{remark}[Length of chains]\label{rem:chain_length}
If $|\im(f)|=m$, every strictly increasing chain of $f$-definable predicates has at most $m+1$ elements, and chains of exactly $m+1$ elements exist.
By the order isomorphism, such a chain is a strictly increasing chain of subsets of an $m$-element set, and each step adds at least one element.
Adding one element at a time from $\emptyset$ to $\im(f)$ gives a chain of length $m+1$.
\end{remark}

\begin{remark}[How ladders can continue]\label{rem:ladders_continue}
Both theorems fix the rules of the game.
An unbounded ladder requires changing them: a packaging map that is not idempotent, or a sequence of interfaces $f_1,f_2,\dots$ in place of a single fixed one.
A changing interface need not have a growing image.
On $X=\mathbb{N}$ the interfaces $f_n=\1_{\{0,\dots,n\}}$ all have a two-point image, yet the predicates $f_1<f_2<\cdots$, each definable through its own interface, form an infinite strictly increasing chain.
Growth becomes necessary when the interfaces are successive refinements, $f_n=g_n\circ f_{n+1}$ for maps $g_n$, as when coordinates are added or the lens is refined.
Then $\Def(f_n)\subseteq\Def(f_{n+1})$, and $g_n$ maps $\im(f_{n+1})$ onto $\im(f_n)$, so $|\im(f_n)|$ is nondecreasing.
If these sizes stay bounded, they are eventually constant; from then on each $g_n$ restricts to a bijection $\im(f_{n+1})\to\im(f_n)$, every predicate definable through $f_{n+1}$ is already definable through $f_n$, and the union of all $\Def(f_n)$ is a finite set.
Along a chain of refinements, an infinite sequence of new distinctions therefore requires $|\im(f_n)|\to\infty$.
In the vocabulary of the introduction, such growth is a form of staging or refinement.
\end{remark}
